\section{Open questions}

The following five questions are grounded in a re-read of the abstract, MeSH headings, and
Kato-lab OA companion papers (PMC10727999 Buglewicz 2023 Cancer Sci.; PMC6467899 Kato 2019 Sci.
Rep.). They flag genuinely open scientific questions --- not merely gaps this backfill could
not close --- and each is paired with concrete next steps that could be pursued with public
funding and modern tooling.

\begin{enumerate}

\item \textbf{At what dose rate does the NHEJ-deficient inverse dose-rate effect (IDRE)
  saturate, and does it invert again in the truly protracted ($<$0.05\,Gy/h, $>$48\,h) regime?}

  \emph{Basis.} The Buglewicz 2024 BBRC paper reports IDRE qualitatively in NHEJ-deficient CHO
  lines but --- based on Kato-lab prior work and typical BBRC experimental scale --- almost
  certainly does not sweep $\dot D$ over more than $\approx$1.5 orders of magnitude, and unlikely
  extends into the truly protracted regime where cells complete multiple divisions during
  exposure and cell-cycle redistribution dominates the survival response. The classic
  Mitchell/Bedford IDRE literature (HeLa, 1979) showed IDRE peaks and then relaxes at very low
  $\dot D$; whether an NHEJ-deficient background follows the same non-monotonic pattern is
  unresolved.

  \emph{Next steps.} Follow-up on the same panel (V3, xrs-5, xrs-6) with a Cs-137 or Co-60 LDR
  facility at $\dot D \in \{0.01, 0.03, 0.1, 0.3, 1.0, 3.0, 10\}$\,Gy/h delivering total doses
  of 2, 4, 8\,Gy. Fit an extended LQ + Lea--Catcheside + IDRE-saturation model
  $\alpha_{\rm eff}(\dot D) = \alpha(1 + \phi\,\dot D_0/(\dot D_0 + \dot D)\,e^{-\dot D/\dot D_{\rm sat}})$
  to locate $\dot D_{\rm sat}$. Simultaneously monitor cell-cycle by EdU/PI flow at hourly bins
  during exposure to disentangle repair-limited IDRE from cell-cycle-redistribution IDRE.

\item \textbf{Does the reduced dose-rate effect in canonical NHEJ-deficient CHO lines translate
  to modern isogenic iPSC-derived knockouts, or is it a CHO-specific artifact of chronic
  aneuploidy and altered checkpoint function in the AA8 lineage?}

  \emph{Basis.} CHO cells (including AA8-derived V3, xrs-5, xrs-6, 51D1) are permanently
  aneuploid with attenuated p53 and altered G1/S checkpoint responses. Every one of the
  $>$40-year-old xrs/irs/V3 lines was selected under acute high-dose-rate stress, potentially
  co-selecting for compensatory phenotypes that shape their LDR response. Whether the same
  DNA-PKcs, Ku80, or XRCC4 loss in a diploid, p53-competent, checkpoint-intact human iPSC
  gives the same IDRE is not established in the Buglewicz 2024 paper or its Kato-lab
  companions.

  \emph{Next steps.} Generate isogenic $PRKDC^{-/-}$, $XRCC6^{-/-}$ (Ku70), $XRCC5^{-/-}$
  (Ku80), and $XRCC4^{-/-}$ knockouts in a diploid iPSC line (WTC-11 or KOLF2.1J) using
  validated CRISPR guides + HDR of a stop cassette. Differentiate into a proliferating
  epithelial or neural-progenitor state to match the CHO cycling context. Repeat
  clonogenic-equivalent (self-renewal or CFU) LDR survival at matched $(D, \dot D)$ grid.
  Directly compare IDRE magnitude and $\dot D_{\rm sat}$ to the CHO panel; publish either
  confirmation (raises translational confidence) or negative result (bounds the CHO-specific
  caveat).

\item \textbf{Is the NHEJ-deficient IDRE additive, synergistic, or epistatic with pharmacologic
  ATM, DNA-PKcs, or PARP inhibition, and does that reveal the rate-limiting step that flips
  dose-rate sparing into dose-rate sensitization?}

  \emph{Basis.} The paper does not test drug combinations. Mechanistically, if IDRE in an
  NHEJ-null background reflects unrepaired DSBs escaping into S/G2 and being processed by
  error-prone HR or alt-EJ, then an ATM inhibitor (KU-60019, AZD1390) should suppress
  downstream damage signaling and possibly attenuate IDRE; a DNA-PKcs inhibitor
  (M3814/peposertib, AZD7648) in a Ku80-null background should be near-epistatic; a PARP
  inhibitor (olaparib, talazoparib) should synergize by blocking alt-EJ backup. Distinguishing
  these three patterns identifies the actual limiting repair step.

  \emph{Next steps.} Repeat LDR clonogenic panel (V3, xrs-5, 51D1, WT) at fixed $\dot
  D\approx0.1$\,Gy/h, 4\,Gy total, with a $3\times4$ drug matrix:
  $\{$DMSO, KU-60019 3\,$\mu$M, AZD7648 1\,$\mu$M, olaparib 1\,$\mu$M$\}\times\{$WT,
  HR$^-$, NHEJ$^-$-1, NHEJ$^-$-2$\}$. Score SF, $\gamma$-H2AX foci at 1, 6, 24\,h, cell-cycle
  at exposure end. Fit isobologram or Bliss/Loewe indices; report which single inhibitor
  reverses or amplifies IDRE most --- that is the rate-limiting node.

\item \textbf{Does the reduced-to-inverse dose-rate effect in cultured NHEJ-deficient CHO cells
  predict in-vivo tumor radioresponse when a tumor with a NHEJ-pathway alteration is treated
  with clinically realistic LDR brachytherapy or protracted external-beam schedules?}

  \emph{Basis.} LDR biology has direct clinical stakes: HDR-vs-LDR brachytherapy decisions
  (prostate, cervical), pulsed-dose-rate schedules, and Lu-177/Ac-225 radioligand therapy all
  deliver at rates comparable to (or slower than) the paper's LDR conditions. If
  NHEJ-compromised tumors show IDRE rather than sparing, LDR modalities become preferentially
  useful for those tumors. Sarcomas, some glioblastomas, and a subset of head-and-neck cancers
  carry $PRKDC$ or Ku alterations, but no in-vivo confirmation of the CHO IDRE finding exists.

  \emph{Next steps.} Two-track approach. Preclinical: xenograft PRKDC-\emph{scid} tumors
  (SCID mouse background is already DNA-PKcs-null) from human $PRKDC$-mutant sarcoma lines vs
  matched WT; treat with an implanted Cs-131 or I-125 LDR source vs matched dose-equivalent
  single-fraction acute; endpoints: tumor growth delay, $\mathrm{TCD}_{50}$, histologic
  mitotic catastrophe. Clinical/epidemiological: mine TCGA + AACR-GENIE for tumors with
  $PRKDC/XRCC4/6/5/LIG4$ loss-of-function; cross-link to clinical radiotherapy records
  (e.g. via institutional cohorts or MSK-IMPACT clinical data) to test whether LDR/pulsed
  schedules give better local control in those tumors.

\item \textbf{Do NHEJ-deficient and HR-deficient CHO lines differ in their LDR response by
  DSB-complexity spectrum --- does IDRE scale with the fraction of clustered/complex DSBs
  produced by high-LET secondaries within the nominally-low-LET $\gamma$-beam, or does
  simple-DSB accumulation alone drive it?}

  \emph{Basis.} $\gamma$-rays produce a spectrum of DSB complexity: predominantly simple
  (single-DSB, single Ku-loading site), with a minor tail ($\approx5$--$15\%$) of complex DSBs
  from local high-LET secondary electrons/$\delta$-rays. NHEJ handles simple DSBs quickly
  ($t_{1/2}\approx20$\,min); complex DSBs require end-resection and slow HR or alt-EJ
  ($t_{1/2}>6$\,h). Whether IDRE in NHEJ-null CHO reflects (a) simple DSBs no longer being
  cleared fast enough, or (b) complex DSBs mis-processed by alt-EJ into lethal chromosomal
  exchanges, is not addressed in the paper (abstract does not mention 53BP1, RIF1, RAD52, or
  complexity readouts).

  \emph{Next steps.} Add a DSB-complexity readout: costain $\gamma$-H2AX / 53BP1 / RAD51 at
  1, 6, 24\,h post-LDR to distinguish simple ($\gamma$-H2AX-only, resolved $<$6\,h) from
  complex (53BP1-persistent $>$24\,h, RAD51-focus-forming) DSBs per line per $\dot D$.
  Complement with pulsed-field gel electrophoresis and a chromosomal-aberration score
  (metaphase spread, dicentric + translocation count). Fit a two-compartment repair model
  (fast simple + slow complex) to check whether IDRE arises from the simple, complex, or both
  compartments in the NHEJ-null background. Cross-check with a proton/carbon-ion LDR arm
  (higher intrinsic complex-DSB fraction) --- Kato-lab has carbon-ion facility access, so this
  is directly actionable.

\end{enumerate}
